\chapterauthor{Lech Tuzinkiewicz}
\chapter[Projektowanie baz \ldots]{Projektowanie baz danych sterowane jakością}
\thispagestyle{empty}
\vspace{-8mm} {\large Lech Tuzinkiewicz} \vspace{8mm}

% chapter_projektowanie.tex


\section{Wprowadzenie}

Obecnie działalność naukowa i gospodarka zdominowane są przez procesy cyfryzacji i~transformacji. Na naszych oczach dokonuje się fundamentalna zmiana paradygmatu danych, przy czym paradygmat należy rozumieć jako koncepcję, perspektywę lub zbiór idei, które wpływają na sposób rozumienia i interpretacji wycinka rzeczywistości w określonej dziedzinie. W tym kontekście dane, które jeszcze kilkanaście lat temu traktowano jako element działalności operacyjnej, obecnie uważane są za produkt o znacznej wartości i mający ogromne znaczenie w kształtowaniu strategii konkurencyjnej. Dane nie są już jedynie zasobem gromadzonym w bazach danych i wykorzystywanym w operacjach biznesowych, ale stanowią zarówno źródło pozyskiwania informacji, jak i istotny zasób majątkowy, czyli walor w postaci aktywu o znaczeniu często przewyższającym tradycyjne aktywa materialne, takie jak środki trwałe czy kapitał finansowy. Aktualnie również uważa się, że dane można postrzegać jako najcenniejszy aktyw we współczesnym świecie.

W erze sztucznej inteligencji i uczenia maszynowego, jakościowo dobre dane mają istotne znaczenie dla osób oraz systemów informatycznych wykorzystywanych w~podejmowaniu decyzji strategicznych w oparciu o analitykę. Znaczenie jakości danych w kontekście ich wartości jest absolutnie fundamentalne. Wyraża się to między innymi powszechnym stwierdzeniem w postaci GIGO, tzn.\ \emph{Garbage In, Garbage Out}, czyli jakość podejmowanych decyzji odpowiada jakości danych, które w danym procesie są wykorzystywane. Podsumowując, jakość danych ma wpływ na decyzje biznesowe oraz modele AI i tym samym mogą być one tylko tak dobre, jak dane, na których się opierają.

Dane gromadzone są w różnych bazach danych (np.\ relacyjnych, NoSQL, dokumentowych itd.), a proces ich projektowania musi uwzględniać aspekt oceny jakości danych, które mają być gromadzone i udostępniane w trakcie realizacji procesów biznesowych.

W dalszej części rozdziału zostanie przedstawiony krótki przegląd literatury obejmujący wybrane aspekty oceny jakościowej procesu projektowania baz danych oraz przykładowe modele jakości mające zastosowanie w tym procesie.

\section{Przegląd literatury}

Projektowanie baz danych jest procesem wieloetapowym, obejmującym tworzenie modeli konceptualnych, logicznych i fizycznych. Jakość tych modeli ma kluczowe znaczenie dla poprawnego działania systemów informatycznych, ponieważ błędy popełnione na etapie projektowania skutkują kosztownymi problemami w implementacji i utrzymaniu. Literatura podkreśla, że ocena jakości powinna być prowadzona iteracyjnie na każdym etapie projektowania, a nie dopiero po wdrożeniu systemu \cite{Dubielewicz2007SPoQE,ISO25012}.

Badania nad jakością modeli danych i baz danych prowadzone są od wielu lat, a ich znaczenie rośnie wraz ze wzrostem roli danych w procesach informatycznych i biznesowych. W literaturze można wyróżnić kilka głównych nurtów badań.

\subsection{Modele jakości danych i ich podstawy normatywne}

Podstawowym punktem odniesienia dla oceny jakości danych są standardy ISO/IEC, w szczególności model jakości danych ISO/IEC 25012:2014 \cite{ISO25012}, który definiuje zestaw charakterystyk jakościowych odnoszących się zarówno do jakości wewnętrznej, jak i zewnętrznej produktu programistycznego. Model ten jest szeroko stosowany w procesach specyfikacji wymagań oraz ewaluacji jakości.

\subsection{Metryki jakości modeli danych}

Gosain w pracy \cite{Gosain2015} przedstawił przegląd metryk jakości modeli danych, wskazując potrzebę ich systematyzacji i dostosowania do praktyki projektowej. Heinrich wraz z zespołem \cite{Heinrich2018} zaproponowali wymagania wobec metryk jakości danych, podkreślając konieczność ich jednoznaczności i powtarzalności. Natomiast Helskyaho i~współautorzy \cite{Helskyaho2024} rozwinęli te zagadnienia w~kontekście Data Vault 2.0, definiując zestaw metryk jakościowych dla tego podejścia.

\subsection{Jakość schematów baz danych}

Sharma i współautorzy \cite{Sharma2018} skupili się na problemie tzw.\ \emph{schema smells}, czyli defektów w~schematach baz danych, które mogą obniżać ich jakość i~utrudniać dalszy rozwój systemu. Podobne zagadnienia były podejmowane w pracach \cite{Kedzia2011,Tuzinkiewicz2013}, gdzie autorzy analizowali kryteria oceny jakości baz danych w kontekście praktycznych zastosowań.

\subsection{Model-Driven Architecture (MDA) jako podstawa projektowania baz danych}

Literatura jednoznacznie wskazuje na zastosowanie podejścia MDA do strukturyzacji procesu projektowania baz danych \cite{Dubielewicz2007Eval,Quad2010,Tuzinkiewicz2013}. MDA dzieli proces wytwarzania oprogramowania na transformacje między modelami:

\begin{itemize}
  \item \textbf{Model konceptualny (PIM -- Platform Independent Model):} Odpowiednik pierwszego etapu, niezależny od technologii. Obejmuje semantyczną zgodność z~dziedziną problemu i poprawność składniową (np.\ zgodność z~notacją ERD/UML).
  \item \textbf{Model logiczny (PSM 1 -- Platform Specific Model 1):} Otrzymywany przez zastosowanie reguł transformacji z modelu PIM. Ten etap wymaga oceny jakości, ponieważ to na nim podejmowane są decyzje rzutujące na wydajność i utrzymanie. Autorzy \cite{Dubielewicz2007Eval} podkreślają, że wybór najlepszego modelu logicznego (z wielu możliwych) powinien opierać się na kryteriach jakościowych, takich jak wydajność czy łatwość utrzymania.
  \item \textbf{Model fizyczny (PSM 2):} Finalny artefakt, gotowy do implementacji w~konkretnym środowisku bazodanowym, który jest produktem końcowym wymagającym oceny jakościowej \cite{Kedzia2011}.
\end{itemize}

Zasady modelowania danych w kontekście MDA oraz identyfikacja kluczowych artefaktów towarzyszących procesowi są szeroko omawiane m.in.  w~\cite{Tuzinkiewicz2013}.

\subsection{Metodyki oceny jakości modeli MDA}

Ważnym nurtem badań są prace dotyczące oceny jakości modeli baz danych w~podejściu MDA. Autorzy w pracy \cite{Dubielewicz2007Eval} opisali opracowaną metodykę QUAD oraz przedstawili przykłady oceny jakości modeli fizycznych powstałych w wyniku transformacji modeli konceptualnych. Wskazali oni, że ocena jakości powinna uwzględniać zarówno wymagania funkcjonalne, jak i niefunkcjonalne.

\subsection{Zarządzanie wymaganiami jakościowymi}

GuerraGarcía i współautorzy \cite{GuerraGarcia2023} zaproponowali metodykę opartą na ISO/IEC 25012 \cite{ISO25012}, która wspiera zarządzanie wymaganiami jakościowymi w~procesach projektowych. Ich podejście łączy perspektywę normatywną z~praktycznymi narzędziami oceny.

\medskip

Podsumowując, analiza literatury wskazuje, że problematyka jakości danych i~modeli danych jest szeroko dyskutowana zarówno w kontekście norm ISO, jak i~praktycznych metryk oraz metod oceny jakości. Istnieje wyraźna tendencja do wykorzystania teoretycznego podejścia (modele jakości, standardy) w praktycznych zastosowaniach. Wnioski z~dotychczasowych badań podkreślają potrzebę dalszej integracji tych podejść, aby zapewnić spójne i~powtarzalne metody oceny jakości baz danych jako produktów końcowych.

\section{Model jakości projektowania baz danych}

Ocena jakości projektowania baz danych łączy perspektywę inżynierii oprogramowania (modele jakości, zgodność ze standardami) z perspektywą bazodanową (normalizacja, integralność, wydajność i utrzymanie). W tym zakresie ma zastosowanie norma jakości SQuaRE (Software Quality Requirements and Evaluation), która jest serią międzynarodowych standardów dotyczących jakości oprogramowania i składa się z sekcji odpowiadających etapom cyklu życia produktu informatycznego:

\begin{itemize}
  \item \textbf{Modele jakości} (np.\ ISO/IEC 25010/25012) -- definiują zbiór charakterystyk i~podcharakterystyk, które opisują kontekst oceny jakościowej. Perspektywa oceny obejmuje:
  \begin{itemize}
    \item \emph{Jakość w użytkowaniu (Quality in Use)} -- ocena z punktu widzenia użytkownika produktu,
    \item \emph{Model jakości (Product Quality Model)} -- opisuje wewnętrzne i~zewnętrzne właściwości ocenianego produktu/artefaktu.
  \end{itemize}
  \item \textbf{Metryki jakości} (np.\ ISO/IEC 25020) -- definiują metryki do pomiaru każdej z~charakterystyk zdefiniowanych w modelu, np.\ „wydajność” poprzez metrykę \emph{czas wykonania zapytania}.
  \item \textbf{Wymagania jakości} (np.\ ISO/IEC 25030) -- przedstawiają zasady definiowania i specyfikowania wymagań jakościowych dla produktu (oprogramowania, bazy danych) na etapie zamawiania i późniejszego rozwoju.
  \item \textbf{Ocena jakości} (np.\ ISO/IEC 25040) -- definiuje proces przeprowadzenia oceny jakości produktu na podstawie zdefiniowanych wymagań i metryk.
  \item \textbf{Dokumentacja oceny} (np.\ ISO/IEC 25041) -- dostarcza szablonów i wymagań dotyczących dokumentowania procesu oceny i jego wyników.
\end{itemize}

Analizując możliwości wykorzystania tej normy do oceny procesu projektowania baz danych, warto przyjrzeć się normie ISO/IEC 25010/25012, która jest podstawowym elementem specyfikacji zakresu oceny jakości produktu, definiując osiem podstawowych charakterystyk:

\begin{itemize}
  \item funkcjonalność (functional suitability),
  \item niezawodność (reliability),
  \item wydajność (performance efficiency),
  \item użyteczność (usability),
  \item bezpieczeństwo (security),
  \item kompatybilność (compatibility),
  \item łatwość utrzymania (maintainability),
  \item przenośność (portability).
\end{itemize}

Normę jakości należy traktować jako zalecenie/podstawę do zdefiniowania sposobu i~zakresu oceny jakości procesu i/lub produktu, wybierając jedynie odpowiednie charakterystyki jakości do rozpatrywanej dziedziny. Opracowując model jakości oceny procesu projektowania baz danych, można wykorzystać strukturę charakterystyk jakości i metody pomiaru, dostosowując je do specyfiki projektowania baz danych.

W tabeli~\ref{tab:iso25000-stages} przedstawiono główne etapy projektowania baz danych wraz z~charakterystykami ISO/IEC 25000 (SQuaRE) oraz przykładowymi miarami i~wagami.


\renewcommand{\arraystretch}{0.9} % giảm chiều cao dòng ~10%
\begin{table}[!ht]
  \centering
  \caption{Charakterystyki ISO/IEC 25000, miary i wagi w kontekście etapów projektowania baz danych}
  \label{tab:iso25000-stages}
  \scriptsize
  \begin{tabularx}{\textwidth}{LLp{5cm}c}
    \toprule
    \textbf{Etap} & \textbf{Charakterystyka (ISO/IEC 25012)} & \textbf{Przykładowa miara} & \textbf{Waga} \\
    \midrule
    Opis rzeczywistości &
    Funkcjonalność &
    \% wymagań zidentyfikowanych &
    0.40 \\
    &
    Niezawodność &
    Liczba sprzecznych reguł biznesowych &
    0.20 \\
    Model konceptualny (UML) &
    Funkcjonalność &
    \% klas odwzorowujących wymagania &
    0.30 \\
    &
    Użyteczność &
    Ocena czytelności diagramu (skala 1--5) &
    0.20 \\
    &
    Niezawodność &
    Liczba konfliktów w modelu &
    0.20 \\
    % \hline
    Model logiczny (SQL) &
    Kompatybilność &
    \% tabel odpowiadających klasom UML &
    0.40 \\
    &
    Łatwość utrzymania &
    Liczba zależności między tabelami &
    0.30 \\
    Implementacja fizyczna (DBMS) &
    Wydajność &
    Średni czas wykonania zapytań (ms) &
    0.30 \\
    &
    Bezpieczeństwo &
    Implementacja mechanizmów kontroli dostępu (np.\ role, uprawnienia) &
    0.20 \\
    Testowanie i walidacja &
    Niezawodność &
    \% przypadków testowych zakończonych sukcesem &
    0.40 \\
    &
    Funkcjonalność &
    Liczba niespełnionych reguł &
    0.30 \\
    &
    Bezpieczeństwo &
    Liczba nieautoryzowanych dostępów &
    -- \\
    Optymalizacja indeksów &
    Wydajność &
    Poprawa czasu odpowiedzi po optymalizacji (\%) &
    0.40 \\
    &
    Łatwość utrzymania &
    Liczba błędów w strukturze indeksów &
    0.30 \\
    \bottomrule
  \end{tabularx}
\end{table}



Analizując propozycję modelu jakości określonego w tabeli~\ref{tab:iso25000-stages}, można zauważyć, że wiele elementów istotnych w kontekście projektowania baz danych nie ma swojego odzwierciedlenia w powiązaniu etapów projektowania oraz przypisanych do nich charakterystyk jakościowych i metryk.

Próba dopasowania modelu jakości do założonego celu wiąże się z modyfikacją zbioru charakterystyk oraz miar. Efekt tego typu działania przedstawia zawartość tabeli~\ref{tab:quality-perspective-mod}. Dodatkowe charakterystyki zostały wyróżnione w opisie i stanowią uzupełnienie poprzedniego modelu.

\begin{table}[!ht]
  \centering
  \caption{Przykładowa modyfikacja perspektywy modelu oceny jakościowej}
  \label{tab:quality-perspective-mod}
  \scriptsize
  \begin{tabularx}{\textwidth}{Lp{6cm}L}
    \toprule
    \textbf{Charakterystyka} & \textbf{Opis} & \textbf{Przykładowa miara} \\
    \midrule
    Dokładność &
    Stopień, w jakim dane poprawnie odzwierciedlają wartości rzeczywiste &
    Wskaźnik błędów w danych (\%) \\
    Kompletność &
    Kompletność atrybutów i ich wartości &
    Procent brakujących wartości \\
    Spójność &
    Brak sprzeczności w danych &
    Liczba niespójnych rekordów \\
    Wiarygodność (źródła) &
    Stopień zaufania do źródła danych &
    Ocena wiarygodności źródła (skala 1--5) \\
    Interpretowalność / użyteczność &
    Łatwość rozumienia i interpretacji danych &
    Ocena interpretowalności (skala 1--5) \\
    Pochodzenie danych (data lineage) &
    Możliwość prześledzenia pochodzenia danych &
    Procent danych z pełnymi metadanymi \\
    Zgodność regulacyjna &
    Zgodność z wymogami prawnymi i regulacyjnymi &
    Liczba naruszeń regulacyjnych \\
    Bezpieczeństwo &
    Ochrona przed nieautoryzowanym dostępem &
    Liczba incydentów bezpieczeństwa \\
    Integralność danych &
    Ochrona integralności i poufności danych &
    Liczba naruszeń integralności \\
    Przenośność &
    Łatwość przenoszenia danych między systemami &
    Czas wymagany na migrację danych \\
    Łatwość utrzymania &
    Możliwość przywrócenia danych po awarii &
    Średni czas odzyskiwania danych \\
    Aktualność &
    Świeżość danych względem czasu użycia &
    Średni interwał odświeżania danych \\
    Wydajność &
    Łatwość uzyskania danych z systemu &
    Średni czas dostępu do danych; czas działania systemu \\
    \bottomrule
  \end{tabularx}
\end{table}

Dokładna analiza modelu jakości przedstawionego w tabeli~\ref{tab:quality-perspective-mod} jest bliższa celowi, ale w dalszym ciągu nie uwzględnia specyfiki projektowania baz danych. Zatem przystępując do opracowania modelu jakości, należy określić perspektywę oceny, tzn.\ przedstawić kluczowe elementy mające wpływ na jakość projektowanego artefaktu. Z punktu widzenia projektowania baz danych możemy rozpatrywać następujące elementy:

\begin{itemize}
  \item \textbf{Jakość opisu wycinka rzeczywistości i wymagań biznesowych}
  \begin{itemize}
    \item Czy opis wycinka rzeczywistości jest kompletny i interpretowalny (zrozumiały) przez uczestników procesu projektowania?
    \item Czy można uzyskać odpowiedzi na pytania „DLACZEGO?” i~„JAKI JEST CEL PRZEDSIĘWZIĘCIA?”.
    \item Czy podstawowe pojęcia biznesowe są poprawnie zdefiniowane w~słowniku (np.\ „Zamówienie”, „Produkt”)?
    \item Czy zdefiniowano klasy obiektów dziedziny wraz z atrybutami, których wartości gwarantują rozróżnialność obiektów?
    \item Czy zdefiniowano reguły biznesowe i ograniczenia dziedzinowe?
    \item Czy uwzględniono planowane modyfikacje modelu danych w bliskiej perspektywie czasowej (w przeciągu 2--3 lat)?
  \end{itemize}
  \item \textbf{Kryteria transformacji konceptualnego modelu danych w środowisku SQL}
  \begin{itemize}
    \item Czy klasy (encje) reprezentują obiekty dziedziny przedmiotowej?
    \item Czy relacje pomiędzy obiektami klas (asocjacje) oraz liczności końców asocjacji wynikają z reguł biznesowych?
    \item Czy klasy asocjacyjne zawierają tylko atrybuty charakteryzujące związek między obiektami?
    \item Czy atrybuty klas wynikają z definicji obiektów wycinka rzeczywistości i~umożliwiają realizację wymagań funkcjonalnych?
  \end{itemize}
  \item \textbf{Kryteria związane z zarządzaniem jakością danych (Data Governance, Quality by Design)}
  \begin{itemize}
    \item Czy definicje i domyślne wartości pól są zgodne z biznesowymi regułami walidacji?
    \item Czy poprawność danych wymuszona jest na poziomie schematu bazy danych poprzez definicję kluczy głównych i obcych, unikalnych indeksów, typów danych, ograniczeń \texttt{CHECK}, \texttt{NOT NULL}?
    \item Czy model zawiera metadane?
    \item Czy dostępność danych zdefiniowana jest na poziomie schematu bazy danych?
    \item Czy uwzględniono mechanizmy kontroli dostępu?
    \item Czy dane wrażliwe są logicznie oddzielone?
    \item Czy określono, które tabele przechowują dane wrażliwe (kluczowe dla bezpieczeństwa i RODO)?
    \item Czy uwzględniono szyfrowanie kolumn lub tabel z danymi poufnymi?
    \item Czy możliwe jest śledzenie, skąd dane pochodzą i jak są przetwarzane (data lineage)?
    \item Czy transakcje gwarantują spójność? Czy ograniczenia integralności zapobiegają pojawianiu się „osieroconych” rekordów?
    \item Czy model danych minimalizuje redundancję poprzez odpowiedni poziom normalizacji i zapobiega powstawaniu niespójności, gdy ta sama informacja jest przechowywana w wielu miejscach?
  \end{itemize}
  \item \textbf{Kryteria związane z wydajnością i dostępnością (Performance by Design)}
  \begin{itemize}
    \item Czy tabele, które będą często łączone, mają zdefiniowane odpowiednie klucze i indeksy?
    \item Czy wybór typów danych jest zgodny z własnościami atrybutów obiektów świata rzeczywistego?
  \end{itemize}
\end{itemize}

\section{Model jakości -- przykład konstrukcji}

Podsumowując wcześniejsze rozważania dotyczące możliwości oceny jakości procesu projektowania bazy danych, można --- bazując na wnioskach z~przeprowadzonej analizy --- zaproponować model porównawczej oceny procesu projektowania baz danych, który łączy kolejne etapy tego procesu, od opisu rzeczywistości do testowania implementacji bazy danych w środowisku SQL. Poniżej zostały określone podstawowe etapy podlegające ocenie jakościowej w procesie projektowania baz danych:

\begin{enumerate}
  \item \textbf{Opis wycinka rzeczywistości (etap analityczny)} \\
  Wejście: narracyjny opis dziedziny lub przypadków użycia. \\
  Wyjście: zbiór pojęć, relacji, reguł biznesowych i ograniczeń dziedzinowych. \\
  Cel: uchwycenie istotnych bytów i zależności.
  \item \textbf{Model konceptualny (np.\ diagram klas UML)} \\
  Wejście: zidentyfikowane pojęcia, relacje, reguły i ograniczenia. \\
  Wyjście: diagram klas UML (z atrybutami, asocjacjami, dziedziczeniem). \\
  Cel: formalizacja pojęć w modelu niezależnym od technologii.
  \item \textbf{Model logiczny i fizyczny (implementacja SQL)} \\
  Wejście: diagram klas lub model logiczny. \\
  Wyjście: schemat relacyjny, tabele, klucze, ograniczenia, widoki. \\
  Cel: odwzorowanie modelu konceptualnego w SQL.
  \item \textbf{Testowanie i walidacja} \\
  Wejście: implementacja + zestaw przypadków testowych (danych). \\
  Wyjście: wyniki testów, ocena zgodności z wymaganiami i semantyką. \\
  Cel: sprawdzenie poprawności, spójności i jakości danych.
\end{enumerate}

Dla każdego etapu zostały określone charakterystyki oraz przykładowe miary jakości (tabela~\ref{tab:etapy-char-miary}).

\begin{table}[!ht]
  \centering
  \caption{Charakterystyki oraz przykładowe miary oceny jakości etapów projektowania}
  \label{tab:etapy-char-miary}
  \scriptsize
  \begin{tabularx}{\textwidth}{LLp{5cm}}
    \toprule
    \textbf{Etap} & \textbf{Charakterystyki} & \textbf{Przykładowe miary} \\
    \midrule
    Opis rzeczywistości &
    Pełność, spójność pojęciowa, jednoznaczność &
    Liczba zidentyfikowanych bytów w stosunku do oczekiwanych; liczba sprzeczności lub niejednoznacznych definicji \\
    % \hline
    Model konceptualny (diagram klas UML) &
    Poprawność semantyczna, złożoność modelu, pokrycie pojęć &
    Liczba klas/asocjacji; wskaźnik pokrycia pojęć (\%); liczba błędnych relacji (np.\ asocjacji bez ról) \\
    % \hline
    Model logiczny (SQL) &
    Zgodność z modelem UML, normalizacja, złożoność implementacji &
    Miary zgodności (mapowanie klas$\rightarrow$tabel, asocjacji$\rightarrow$FK); poziom normalizacji (np.\ \% tabel w 3NF); liczba obiektów SQL \\
    % \hline
    Testowanie &
    Zgodność wyników z oczekiwaniami, wydajność, spójność danych &
    Liczba błędnych wyników; czas uzyskania danych z bazy danych \\
    \bottomrule
  \end{tabularx}
\end{table}

Zaproponowane charakterystyki stanowią zarys modelu oceny jakości procesu projektowania baz danych. Ocena końcowa jakości $Q$ jest wyrażona za pomocą funkcji agregującej oceny cząstkowe dla każdego etapu:
\begin{equation}
  Q = f(Q_{\mathrm{opis}}, Q_{\mathrm{UML}}, Q_{\mathrm{SQL}}, Q_{\mathrm{test}}),
\end{equation}
gdzie:
\begin{itemize}
  \item $Q_{\mathrm{opis}}$ -- jakość opisu rzeczywistości,
  \item $Q_{\mathrm{UML}}$ -- jakość modelu UML,
  \item $Q_{\mathrm{SQL}}$ -- jakość schematu SQL,
  \item $Q_{\mathrm{test}}$ -- jakość testowania,
  \item wagi etapów spełniają warunek: suma wag $=1$.
\end{itemize}

Dla każdego etapu procesu projektowania należy przeprowadzić jakościową ocenę tworzonych artefaktów, biorąc pod uwagę charakterystyki, miary i sposób ich interpretacji.

\subsection*{Etap 1: Opis wycinka rzeczywistości}

\subsubsection*{Cel etapu}

Przygotowanie zrozumiałego i kompletnego opisu dziedziny problemowej, który stanowi punkt wyjścia do modelowania konceptualnego danych i obejmuje:
\begin{itemize}
  \item identyfikację bytów świata rzeczywistego (encji),
  \item ich cechy (atrybuty),
  \item relacje i ograniczenia dziedzinowe,
  \item reguły biznesowe określające dopuszczalne stany i operacje.
\end{itemize}

Opis może mieć formę:
\begin{itemize}
  \item narracyjną (tekstową),
  \item zestawu przypadków użycia,
  \item dokumentacji wymagań,
  \item scenariuszy procesów biznesowych.
\end{itemize}

Główne charakterystyki jakości opisu rzeczywistości zostały przedstawione w~tabeli~\ref{tab:opis-char}.

\begin{table}[!ht]
  \centering
  \caption{Charakterystyki jakości opisu rzeczywistości}
  \label{tab:opis-char}
  \scriptsize
  \begin{tabularx}{\textwidth}{cLLp{5cm}}
    \toprule
    \textbf{Nr} & \textbf{Charakterystyka} & \textbf{Znaczenie} & \textbf{Przykład} \\
    \midrule
    O1 &
    Pełność bytów i atrybutów &
    Czy opis identyfikuje wszystkie istotne obiekty oraz ich cechy. &
    Opis zawiera byty: Klient, Zamówienie, Produkt oraz ich atrybuty, np.\ „nazwa”, „cena”. \\
    % \hline
    O2 &
    Identyfikowalność bytów &
    Czy każdy byt ma naturalny identyfikatór jednoznacznie rozróżniający obiekty. &
    Klient: PESEL lub e-mail; Produkt: kod EAN. \\
    % \hline
    O3 &
    Reguły biznesowe &
    Czy opis zawiera zdefiniowane reguły określające poprawne stany, relacje i ograniczenia. &
    „Każde zamówienie musi mieć co najmniej jeden produkt”; „Cena nie może być ujemna”. \\
    % \hline
    O4 &
    Ograniczenia dziedzinowe &
    Czy wartości atrybutów są ograniczone zgodnie z dziedziną (zakresy, typy, dopuszczalne wartości). &
    Wiek 0--120, status: \{„aktywne”, „zawieszone”\}. \\
    % \hline
    O5 &
    Spójność i jednoznaczność &
    Czy opis nie zawiera sprzecznych lub niejasnych reguł i terminów. &
    Brak sprzecznych założeń typu: „Klient może mieć jedno zamówienie” vs „Klient może mieć wiele zamówień”. \\
    % \hline
    O6 &
    Strukturalność i czytelność &
    Czy opis ma uporządkowaną formę i jest zrozumiały. &
    Użyto sekcji, glosariusza, przypadków użycia. \\
    \bottomrule
  \end{tabularx}
\end{table}





W tabeli~\ref{tab:opis-metryki} zostały zaproponowane miary jakości dotyczące etapu charakterystyki dziedziny przedmiotowej.

\begin{table}[!ht]
  \centering
  \caption{Przykładowe miary jakości (metryki) dla etapu opisu rzeczywistości}
  \label{tab:opis-metryki}
  \scriptsize
  \begin{tabularx}{\textwidth}{cLp{5cm}C}
    \toprule
    \textbf{Kod} & \textbf{Miara} & \textbf{Opis pomiaru} & \textbf{Zakres / interpretacja} \\
    \midrule
    M1 &
    Wskaźnik pełności bytów &
    Liczba bytów zidentyfikowanych / liczba bytów oczekiwanych &
    0--1 \\
    % \hline
    M2 &
    Wskaźnik identyfikowalności &
    Liczba bytów z określonym naturalnym identyfikatorem / liczba wszystkich bytów &
    0--1 \\
    % \hline
    M3 &
    Wskaźnik pokrycia reguł biznesowych &
    Liczba reguł opisanych / liczba reguł oczekiwanych (np.\ na podstawie specyfikacji domeny) &
    0--1 \\
    % \hline
    M4 &
    Wskaźnik określenia ograniczeń dziedzinowych &
    Liczba atrybutów z określoną dziedziną wartości / liczba wszystkich atrybutów &
    0--1 \\
    % \hline
    M5 &
    Wskaźnik spójności semantycznej &
    Liczba sprzeczności lub niejednoznaczności / liczba elementów (przeliczany na 1$-\frac{\text{błędy}}{\text{liczba elementów}}$) &
    0--1 \\
    % \hline
    M6 &
    Zrozumiałość i strukturalność &
    Ocena ekspercka (skala 1--5) lub automatyczna analiza struktury (znormalizowana) &
    0--1 \\
    \bottomrule
  \end{tabularx}
\end{table}


Funkcja oceny jakości etapu reprezentowana jest w postaci wyrażenia:
\begin{equation}
  Q_{\mathrm{opis}} = g(w_1 M1, w_2 M2, \dots, w_6 M6),
\end{equation}
gdzie $w_i$ stanowią propozycję wartości wag poszczególnych miar (np.\ $w_1 = \dots = w_6 = \frac{1}{6}$), a wynik $Q_{\mathrm{opis}}$ reprezentuje jakość semantyczną opisu dziedziny.

Wartość funkcji oceny można interpretować w kontekście przykładowych wartości progowych:
\begin{itemize}
  \item $Q_{\mathrm{opis}} \ge 0.85$: opis kompletny, jednoznaczny, dobrze ustrukturyzowany;
  \item $0.70 \le Q_{\mathrm{opis}} < 0.85$: część pojęć brak lub niejasne zależności;
  \item $Q_{\mathrm{opis}} < 0.70$: opis wymaga poważnej rekonstrukcji przed modelowaniem UML.
\end{itemize}

\subsection*{Etap 2: Model konceptualny (diagram klas UML/ERD)}

\subsubsection*{Cel etapu}

Zbudowanie formalnego, spójnego i wystarczająco szczegółowego modelu pojęciowego, który wiernie odwzorowuje rzeczywistość, a zarazem może być podstawą zdefiniowania modelu logicznego bazy danych SQL. Model powinien zagwarantować:
\begin{itemize}
  \item poprawne odwzorowanie pojęć i reguł z etapu 1,
  \item zawieranie klas z określonymi atrybutami i naturalnymi identyfikatorami,
  \item uwzględnienie ograniczeń dziedzinowych i reguł biznesowych w strukturze diagramu klas UML (asocjacje, liczności, OCL itp.).
\end{itemize}

Wybrane charakterystyki zostały przedstawione w tabeli~\ref{tab:uml-char}.

\begin{table}[!ht]
  \centering
  \caption{Główne charakterystyki w kontekście oceny jakości modelu UML}
  \label{tab:uml-char}
  \scriptsize
  \begin{tabularx}{\textwidth}{cLp{4cm}p{4cm}}
    \toprule
    \textbf{Nr} & \textbf{Charakterystyka} & \textbf{Znaczenie} & \textbf{Przykład} \\
    \midrule
    C1 &
    Poprawność semantyczna (semantic correctness) &
    Zgodność modelu z opisem rzeczywistości -- czy klasy, atrybuty i relacje odpowiadają rzeczywistym pojęciom i powiązaniom. &
    Czy zidentyfikowane klasy (np.\ Klient, Zamówienie) odpowiadają pojęciom z opisu dziedziny? \\
    C2 &
    Poprawność odwzorowania reguł biznesowych &
    Reguły opisane w etapie 1 są poprawnie odwzorowane w UML (np.\ poprzez kardynalności, ograniczenia, OCL). &
    Reguła „Każde zamówienie ma co najmniej jeden produkt” odwzorowana jako relacja 1..* między Zamówieniem a Produktem. \\
    C3 &
    Złożoność modelu (complexity) &
    Stopień skomplikowania struktury modelu (liczba klas, relacji, poziomy dziedziczenia). &
    Czy model jest wystarczająco szczegółowy, ale nie przeładowany? \\
    C4 &
    Zgodność z ograniczeniami dziedzinowymi &
    Określono typy i zakresy atrybutów zgodnie z opisem dziedziny. &
    Atrybut „płeć: znak [K, M]”. \\
    C5 &
    Pokrycie pojęć (concept coverage) &
    W jakim stopniu model UML odwzorowuje pojęcia i relacje z opisu rzeczywistości. &
    Czy wszystkie kluczowe elementy z opisu zostały ujęte w modelu? \\
    C6 &
    Czytelność i przejrzystość (readability) &
    Łatwość interpretacji modelu przez człowieka. &
    Czy diagram jest logicznie pogrupowany, a nazwy są jednoznaczne? \\
    C7 &
    Identyfikowalność obiektów &
    Każda klasa posiada naturalne atrybuty jednoznacznie identyfikujące obiekty klas. &
    Klient $\rightarrow$ PESEL; Zamówienie $\rightarrow$ (numer zamówienia, rok). \\
    \bottomrule
  \end{tabularx}
\end{table}



Przykładowe miary przedstawiono w tabeli~\ref{tab:uml-metryki}.

\begin{table}[!ht]
  \centering
  \caption{Przykładowe miary jakości (metryki) dla modelu UML}
  \label{tab:uml-metryki}
  \scriptsize
  \begin{tabularx}{\textwidth}{cLp{4.5cm}p{3.5cm}}
    \toprule
    \textbf{Kod} & \textbf{Miara} & \textbf{Opis sposobu pomiaru} & \textbf{Zakres / interpretacja} \\
    \midrule
    M1 &
    Wskaźnik zgodności semantycznej &
    Liczba pojęć poprawnie odwzorowanych / liczba pojęć z opisu rzeczywistości &
    0--1 \\
    % \hline
    M2 &
    Pokrycie reguł biznesowych &
    Liczba reguł z etapu 1 odwzorowanych w UML (asocjacje, ograniczenia, OCL) / liczba wszystkich reguł &
    0--1 (po przeliczeniu na „poprawność”) \\
    % \hline
    M3 &
    Miara złożoności modelu &
    Liczba klas + liczba asocjacji + liczba atrybutów (lub złożona miara, np.\ Halsteada) &
    Znormalizowana do [0,1] (np.\ optymalny zakres 20--40 klas) \\
    % \hline
    M4 &
    Wskaźnik spójności wewnętrznej &
    Liczba niespójnych relacji / liczba wszystkich relacji (przeliczany na 1$-\frac{\text{niespójności}}{\text{liczba relacji}}$) &
    0--1 \\
    % \hline
    M5 &
    Pokrycie pojęć &
    Liczba pojęć z opisu rzeczywistości odwzorowanych w UML / liczba pojęć z opisu &
    0--1 \\
    % \hline
    M6 &
    Czytelność modelu &
    Ocena ekspercka (skala 1--5) lub metryka automatyczna (np.\ stosunek etykiet tekstowych do linii) &
    0--1 \\
    % \hline
    M7 &
    Wskaźnik identyfikowalności klas &
    Liczba klas z naturalnym identyfikatorem / liczba wszystkich klas &
    0--1 \\
    \bottomrule
  \end{tabularx}
\end{table}


Funkcja oceny jakości konceptualnego modelu danych $Q_{\mathrm{UML}}$ może być przedstawiona jako:
\begin{equation}
  Q_{\mathrm{UML}} = h(w_1 M1, w_2 M2, \dots, w_7 M7),
\end{equation}
gdzie przykładowo $w_1, w_2, w_5$ mają większe znaczenie (np.\ $w_1 = w_2 = w_5 = 0{.}18$), a~pozostałe są niższe (np.\ $w_3 = w_4 = w_6 = w_7 = 0{.}115$), przy zachowaniu sumy wag równej 1.

Przykładowa interpretacja jakości konceptualnego modelu danych przedstawiona jest w tabeli~\ref{tab:uml-ocena}.

\begin{table}[ht]
  \centering
  \caption{Interpretacja jakości konceptualnego modelu danych}
  \label{tab:uml-ocena}
  \scriptsize
  \begin{tabularx}{\textwidth}{LLp{5cm}}
    \toprule
    \textbf{Zakres} & \textbf{Ocena jakości} & \textbf{Charakterystyka} \\
    \midrule
    0.85--1.00 & Wysoka jakość & Model pełny, poprawny, czytelny \\
    0.70--0.85 & Dobra jakość & Drobne braki, możliwe uproszczenia \\
    0.50--0.70 & Średnia jakość & Częściowa niezgodność lub złożoność \\
    $< 0.50$ & Niska jakość & Model wymaga przebudowy \\
    \bottomrule
  \end{tabularx}
\end{table}

Etap 2 można dodatkowo zweryfikować względem etapu 1, np.\ biorąc pod uwagę następujące kryteria:
\begin{itemize}
  \item każda klasa UML $\leftrightarrow$ byt z etapu 1,
  \item każdy atrybut UML $\leftrightarrow$ cecha z etapu 1,
  \item każda asocjacja $\leftrightarrow$ relacja lub reguła biznesowa,
  \item każda reguła OCL $\leftrightarrow$ ograniczenie dziedzinowe.
\end{itemize}

Wymienione kryteria pozwalają zdefiniować wskaźnik zgodności między etapami $Q_{\mathrm{zgod}}$.

\subsection*{Etap 3: Model logiczny / fizyczny (SQL)}

Zakładając akceptowalny poziom jakości konceptualnego modelu danych, można przystąpić do opracowania modelu logicznego i fizycznego danych (w tym przypadku modelu relacyjnego reprezentowanego na poziomie SQL). W modelu należy zaimplementować reguły biznesowe, ograniczenia dziedzinowe i elementy wyspecyfikowane w poprzednich etapach oraz zdefiniować mierzalne kryteria oceny jakości implementacji.

\subsubsection*{Cel etapu}

Odwzorowanie konceptualnego modelu danych (diagram klas UML) w poprawny, spójny i wydajny schemat relacyjny zawierający definicje tabel wraz z ograniczeniami, indeksy, widoki danych, zapytania SQL, procedury, zapewniając:
\begin{itemize}
  \item identyfikowalność bytów (klucze główne z naturalnych identyfikatorów, gdy to możliwe),
  \item implementację reguł biznesowych i ograniczeń dziedzinowych (\texttt{CHECK}, \texttt{FK}, \texttt{UNIQUE}, triggery, procedury),
  \item zgodność z wymaganiami wydajnościowymi i operacyjnymi.
\end{itemize}

Ocenę jakości modelu logicznego/fizycznego można zrealizować, wykorzystując charakterystyki jakości określone w tabeli~\ref{tab:sql-char}.

\begin{table}[!ht]
  \centering
  \caption{Podstawowe charakterystyki jakości implementacji w SQL}
  \label{tab:sql-char}
  \scriptsize
  \begin{tabularx}{\textwidth}{cLp{7cm}}
    \toprule
    \textbf{Nr} & \textbf{Charakterystyka} & \textbf{Znaczenie} \\
    \midrule
    S1 & Zgodność semantyczna z UML & Schemat SQL wiernie odwzorowuje klasy, atrybuty, relacje, identyfikatory i reguły biznesowe. \\
    S2 & Integralność referencyjna & Wszystkie relacje z UML są reprezentowane przez klucze obce lub tabele pośrednie. \\
    S3 & Egzekwowanie ograniczeń dziedzinowych i reguł biznesowych &
    \texttt{CHECK}, \texttt{NOT NULL}, \texttt{UNIQUE}, triggery implementują zdefiniowane ograniczenia. \\
    S4 & Poziom normalizacji / redukcja redundancji &
    Tabele są właściwie znormalizowane (zwykle 3NF) lub celowo zdenormalizowane w uzasadnionych przypadkach. \\
    S5 & Wydajność i skalowalność &
    Schemat i indeksy umożliwiają akceptowalną szybkość zapytań i przetwarzania. \\
    S6 & Bezpieczeństwo i separacja uprawnień &
    Odpowiednie role, uprawnienia, maskowanie danych. \\
    S7 & Utrzymywalność i rozszerzalność &
    Schemat jest czytelny, udokumentowany i łatwy do modyfikacji (zmiana wymagań, migracje). \\
    S8 & Odporność na błędy i transakcyjność &
    Mechanizmy transakcyjne, ograniczenia, backup/restore. \\
    S9 & Testowalność &
    Zdefiniowane testy (jednostkowe, integracyjne, dane) sprawdzające reguły i jakość danych. \\
    \bottomrule
  \end{tabularx}
\end{table}

Zaproponowane poniżej miary jakości (metryki) zostały podzielone na metryki strukturalne (statyczne) i dynamiczne (wynikające z testów/pomiarów). Każda metryka jest znormalizowana do [0,1].

\paragraph{Metryki strukturalne}

\begin{itemize}
  \item M1 -- wskaźnik poprawnego mapowania: liczba elementów poziomu poprawnie odwzorowanych w SQL (klasy$\rightarrow$tabele, atrybuty$\rightarrow$kolumny, identyfi\-katory$\rightarrow$PK/\texttt{UNIQUE}, relacje$\rightarrow$FK/tabele asocjacyjne) / liczba elementów UML.
  \item M2 -- pokrycie reguł biznesowych: liczba reguł biznesowych i ograniczeń dziedzinowych zrealizowanych w schemacie (\texttt{CHECK}/\texttt{FK}/triggery/procedury) / liczba reguł z etapu 1.
  \item M3 -- pokrycie ograniczeń dziedzinowych: liczba atrybutów z zadeklarowanymi typami i ograniczeniami / liczba atrybutów.
  \item M4 -- stopień normalizacji: indeks opisujący, ile tabel spełnia 3NF (lub inny docelowy poziom normalizacji).
  \item M5 -- unikalność i identyfikowalność: odsetek klas, dla których naturalny identyfikator z etapu 1.\ został odwzorowany jako PK lub \texttt{UNIQUE}.
\end{itemize}

\paragraph{Metryki dynamiczne / testowe}

\begin{itemize}
  \item M6 -- wskaźnik integralności danych: $1 - \frac{\text{liczba naruszeń integralności}}{\text{liczba testów integralności}}$.
  \item M7 -- wskaźnik sukcesu testów funkcjonalnych: liczba testów funkcjonalnych zakończonych sukcesem / liczba wszystkich testów.
  \item M8 -- wydajność zapytań: znormalizowana metryka złożona (średni czas zapytania, 95.\ percentyl, przepustowość).
  \item M9 -- wskaźnik jakości danych: $1 - \frac{\text{liczba rekordów niezgodnych z ograniczeniami}}{\text{liczba sprawdzonych rekordów}}$.
  \item M10 -- czas wykonania migracji/wdrożenia: znormalizowana miara (niższy czas $\rightarrow$ wyższa ocena).
  \item M11 -- pokrycie testów: \% przypadków reguł biznesowych pokrytych testami automatycznymi.
\end{itemize}

Przykładowe wartości wag dla metryk strukturalnych i dynamicznych można dobrać tak, aby preferować poprawność odwzorowania i implementację ograniczeń.

\subsection*{Etap 4: Testowanie i walidacja modelu danych (SQL)}

Na etapie testowania brane są pod uwagę następujące rodzaje testów:
\begin{itemize}
  \item testy jednostkowe schematu -- sprawdzenie ograniczeń kolumn (typy, \texttt{NOT NULL}, \texttt{UNIQUE}, \texttt{CHECK}),
  \item testy integracyjne -- tworzenie/aktualizacja rekordów zgodnie ze scenariuszami (sprawdzanie \texttt{FK}, reguł, transakcji),
  \item testy regresyjne reguł biznesowych -- sprawdzanie wszystkich reguł z~etapu 1,
  \item testy wydajnościowe -- pomiar latencji, przepustowości, 95.\ percentyla,
  \item testy jakości danych (data quality) -- skrypty sprawdzające zgodność typów, zakresów, duplikaty,
  \item testy migracyjne -- odtwarzanie migracji na podstawie kopii danych testowych,
  \item testy bezpieczeństwa i uprawnień -- role, \texttt{GRANT}/\texttt{REVOKE}, widoczność danych,
  \item testy regresyjne/automatyczne w ramach CI/CD.
\end{itemize}

Na podstawie metryk M6--M11 dla każdej kategorii testów można zdefiniować wskaźniki oraz miary umożliwiające określenie poziomu jakości etapu 4 (tabela~\ref{tab:testy-metryki}).

\begin{table}[ht]
  \centering
  \caption{Wskaźniki oraz miary jakości etapu testów}
  \label{tab:testy-metryki}
  \scriptsize
  \begin{tabularx}{\textwidth}{cLp{4.5cm}L}
    \toprule
    \textbf{Nr} & \textbf{Wskaźnik} & \textbf{Miara} & \textbf{Komentarz} \\
    \midrule
    T1 &
    Wskaźnik zgodności strukturalnej &
    Odsetek poprawnych typów, kluczy, relacji &
    Poprawność typów, kluczy, \texttt{FK}, nazw \\
    T2 &
    Pokrycie reguł biznesowych &
    Odsetek reguł z etapu 1 egzekwowanych w SQL &
    Ile reguł z etapu 1 działa w SQL \\
    T3 &
    Integralność danych &
    Odsetek danych spełniających ograniczenia &
    Naruszenia PK, \texttt{FK}, zakresów, obowiązkowości atrybutów \\
    T4 &
    Wydajność &
    Znormalizowany czas wykonania zapytań &
    Im mniejszy czas, tym wyższa ocena (do 1.0) \\
    T5 &
    Jakość danych (DQ) &
    Procent poprawnych rekordów &
    Jakość danych testowych/importowanych \\
    T6 &
    Migracja &
    Czas wdrożenia względem limitu &
    Walidacja procesu wdrożenia i migracji \\
    T7 &
    Bezpieczeństwo &
    Odsetek poprawnych reguł dostępu i uprawnień &
    Role, \texttt{GRANT}/\texttt{REVOKE}, maskowanie \\
    T8 &
    Regresja / automatyzacja &
    Pokrycie testów automatycznych &
    CI/CD -- stabilność w czasie \\
    \bottomrule
  \end{tabularx}
\end{table}


Funkcja oceny jakości etapu testowania może być zapisana jako:
\begin{equation}
  Q_{\mathrm{test}} = k(v_1 T1, v_2 T2, \dots, v_8 T8),
\end{equation}
gdzie $v_i$ to wagi przypisane poszczególnym wskaźnikom (np.\ większa waga dla T2, T3, T4, T5).

\subsection*{Łączna ocena całego procesu projektowego}

Po uzyskaniu ocen cząstkowych z etapów 1--4 możemy obliczyć ocenę końcową projektu:
\begin{equation}
  Q_{\mathrm{projekt}} = F(\alpha_1 Q_{\mathrm{opis}}, \alpha_2 Q_{\mathrm{UML}}, \alpha_3 Q_{\mathrm{SQL}}, \alpha_4 Q_{\mathrm{test}}),
\end{equation}
gdzie:
\begin{itemize}
  \item $Q_{\mathrm{opis}}$ -- jakość charakterystyki dziedziny przedmiotowej,
  \item $Q_{\mathrm{UML}}$ -- jakość konceptualnego modelu danych,
  \item $Q_{\mathrm{SQL}}$ -- jakość implementacji w środowisku SQL,
  \item $Q_{\mathrm{test}}$ -- jakość etapu testowania i walidacji,
  \item $\alpha_i$ -- wagi etapów (suma wag $=1$).
\end{itemize}

\section{Podsumowanie i wnioski}

Projektowanie baz danych zapewniające uzyskanie dobrej jakości bazy danych jest poważnym problemem. Przedstawiona w pracy metoda analizy i~oceny jakościowej procesu i produktu, bazująca na normie jakości ISO/IEC 25012:2014, daje szansę weryfikowania i oceny artefaktów na wszystkich etapach procesu projektowania. Wymaga to jednak konsekwentnego przestrzegania zdefiniowanych zasad.

Poprawny jakościowo proces projektowania baz danych powinien zapewnić pozytywne odpowiedzi na następujące pytania:
\begin{itemize}
  \item Czy model danych jest zrozumiały dla biznesu i czy bezpośrednio wspiera jego cele?
  \item Czy jakość danych jest uwzględniona w schemacie bazy danych poprzez implementację ograniczeń i reguł biznesowych?
  \item Czy schemat bazy danych na poziomie definicyjnym uwzględnia integralność i~poufność danych?
  \item Czy uwzględniony jest aspekt wydajności na poziomie schematu bazy danych w~kontekście zbioru operacji na danych?
  \item Czy możliwe jest zarządzanie metadanymi?
  \item Czy projekt bazy danych jest przygotowany na potencjalne zmiany?
\end{itemize}

Ponadto istotne jest wykonywanie iteracyjnej oceny jakości na każdym etapie projektowania (konceptualnym, logicznym, fizycznym), a nie tylko po wdrożeniu. Ramy MDA oferują strukturę transformacji, a normy ISO zapewniają ustandaryzowane wymiary jakości. Główna luka badawcza leży w standaryzacji metryk i~procesu walidacji modeli, co jest niezbędne, aby obiektywna ocena mogła stanowić podstawę do wyboru optymalnej ścieżki transformacji (np.\ z PIM do PSM).

\section{Kierunki dalszych badań i prac}

Pierwszym kierunkiem badań jest opracowanie zintegrowanego, wielowarstwowego modelu jakości, obejmującego zarówno modele konceptualne, jak i logiczne oraz fizyczne. Model ten powinien być zgodny z ISO/IEC 25012 oraz jednocześnie powiązany z procesami transformacji MDA, umożliwiając śledzenie jakości w sposób ciągły i powtarzalny.

Drugim kierunkiem jest budowa uniwersalnego zestawu metryk, możliwego do stosowania niezależnie od używanej notacji (ERD, UML, Data Vault) czy technologii docelowej. Zestaw taki powinien obejmować zarówno metryki strukturalne, semantyczne, jak i~transformacyjne, a także umożliwiać adaptację do specyfiki różnych domen danych.

Kolejnym kierunkiem rozwoju jest automatyzacja oceny jakości, obejmująca narzędzia CASE zdolne do:
\begin{itemize}
  \item automatycznej detekcji \emph{schema smells},
  \item szacowania jakości transformacji PIM$\rightarrow$PSM,
  \item monitorowania zgodności modeli z wymaganiami ISO,
  \item generowania raportów jakościowych na podstawie metryk.
\end{itemize}

Automatyzacja ta mogłaby stworzyć podstawę do pełnej integracji jakości z~procesem DevOps/DataOps, co jest kierunkiem szczególnie istotnym w systemach o wysokiej dynamice zmian (np.\ hurtownie danych).

Czwartym kierunkiem jest rozwój metryk empirycznych, weryfikowanych w~dużych, przemysłowych środowiskach danych. Badania takie umożliwiłyby weryfikację praktycznej przydatności zaproponowanych metryk oraz ich wpływu na jakość wdrażanych systemów.

Piątym kierunkiem jest dalszy rozwój badań nad \emph{smell-based metrics}, czyli metrykami wyprowadzonymi z identyfikowanych wzorców błędów w~schematach. W połączeniu z uczeniem maszynowym mogłoby to umożliwić predykcyjną ocenę jakości modeli danych na podstawie ich struktury oraz historii zmian.

Ostatni kierunek badań dotyczy kwalifikacji jakości danych w kontekście jakości modeli. Potrzebne są badania, które umożliwią powiązanie jakości struktury danych (modelu) z jakością instancji danych operacyjnych i analitycznych. Może to doprowadzić do powstania pełnego, holistycznego modelu jakości obejmującego zarówno artefakty projektowe, jak i dane produkcyjne.


% =====================
% References
% =====================
% \section*{References}
% \addcontentsline{toc}{section}{References}

\begin{thebibliography}{99}

\bibitem{Dubielewicz2007SPoQE}
Dubielewicz, I., Hnatkowska, B., Huzar, Z., Tuzinkiewicz, L.:
An approach to software quality specification and evaluation (SPoQE).
In: IFIP International Federation for Information Processing, vol. 227,
pp. 155--166 (2007)

\bibitem{Dubielewicz2007Eval}
Dubielewicz, I., Hnatkowska, B., Huzar, Z., Tuzinkiewicz, L.:
Evaluation of MDA/PSM database model quality in the context of selected non-functional requirements.
In: Second International Conference on Dependability of Computer Systems (DepCoS-RELCOMEX 2007),
Szklarska Poręba, Polska, pp. 19--26 (2007)

\bibitem{Quad2010}
Dubielewicz, I., Hnatkowska, B., Huzar, Z., Tuzinkiewicz, L.:
Metodyka QUAD. Sterowanie jakością, wytwarzanie aplikacji bazodanowych.
Oficyna Wydawnicza Politechniki Wrocławskiej, Wrocław (2010)

\bibitem{Gosain2015}
Gosain, A.:
Literature review of data model quality metrics.
Procedia Computer Science 48, 236--243 (2015).
https://doi.org/10.1016/j.procs.2015.04.056

\bibitem{GuerraGarcia2023}
GuerraGarcía, C., Nikiforova, A., Jiménez, S., PérezGonzalez, H., RamírezTorres, M., OntañonGarcía, L.:
ISO/IEC 25012-based methodology for managing data quality requirements.
Information Systems (2023).
https://doi.org/10.1016/j.is.2023.101123

\bibitem{Heinrich2018}
Heinrich, B., Hristova, D., Klier, M., Schiller, A., Szubartowicz, M.:
Requirements for data quality metrics.
Journal of Data and Information Quality 9(2), 1--32 (2018).
https://doi.org/10.1145/3148238

\bibitem{Helskyaho2024}
Helskyaho, H., Ruotsalainen, L., Männistö, T.:
Defining data model quality metrics for Data Vault 2.0.
Inventions 9(1), 21 (2024).
https://doi.org/10.3390/inventions9010021

\bibitem{ISO25012}
ISO/IEC 25012:2014 Software engineering --- Software product Quality Requirements and Evaluation (SQuaRE) --- Data quality model.
Edition 1, confirmed current 2025.
International Organization for Standardization (2014)

\bibitem{Kedzia2011}
Kędzia, P., Tuzinkiewicz, L.:
Ocena jakości baz danych.
Studia Informatica 32(2A), 131--142 (2011)

\bibitem{Sharma2018}
Sharma, T., Fragkoulis, M., Rizou, S., Bruntink, M., Spinellis, D.:
Smelly relations: Measuring and understanding database schema quality.
In: Proc. 40th International Conference on Software Engineering (ICSE 2018) (2018).
https://doi.org/10.1145/3183519.3183540

\bibitem{Tuzinkiewicz2013}
Tuzinkiewicz, L.:
Ocena jakości baz danych.
Computer Science (2013)



\end{thebibliography}


